\chapter*{शब्दावली (Glossary)}
\addcontentsline{toc}{chapter}{शब्दावली (Glossary)}
\markboth{शब्दावली}{शब्दावली}

इस पुस्तक में प्रयुक्त प्रमुख वैज्ञानिक एवं तकनीकी शब्दों का संक्षिप्त अर्थ नीचे
अंग्रेज़ी वर्णक्रम (alphabetical order) में दिया गया है।

\begin{longtable}{@{}p{5.2cm} p{8.3cm}@{}}
\toprule
\rowcolor{cvBlue}\hdc{शब्द (Term)} & \hdc{अर्थ (Meaning)} \\
\midrule
\endfirsthead
\toprule
\rowcolor{cvBlue}\hdc{शब्द (Term)} & \hdc{अर्थ (Meaning)} \\
\midrule
\endhead
\bottomrule
\multicolumn{2}{r}{\small\itshape अगले पृष्ठ पर जारी \dots} \\
\endfoot
\bottomrule
\endlastfoot

\textbf{Accountability} (जवाबदेही) & किसी AI-निर्णय के परिणाम के लिए उत्तरदायी मनुष्य का निर्धारण। \\
\textbf{Accuracy} (शुद्धता) & माप या पूर्वानुमान का सही मान के कितना निकट होना। \\
\textbf{AI} (कृत्रिम बुद्धिमत्ता) & मशीनों द्वारा मानव-जैसे बौद्धिक कार्य करना। देखें Artificial Intelligence। \\
\textbf{AlphaFold} & DeepMind की वह AI जिसने प्रोटीन की 3D संरचना का पूर्वानुमान किया (2021)। \\
\textbf{Analysis} (विश्लेषण) & ग्राफ़ एवं सांख्यिकी द्वारा data में pattern देखना (डेटा-विज्ञान का चरण~3)। \\
\textbf{Artificial Intelligence} (कृत्रिम बुद्धिमत्ता) & मशीनों द्वारा देखना, सीखना, तर्क एवं निर्णय जैसे बौद्धिक कार्य करना। \\
\textbf{Bases} (क्षारक) & DNA के चार “अक्षर” - A, T, G, C। \\
\textbf{Bias} (पक्षपात) & एकपक्षीय data के कारण AI के झुके हुए या अनुचित परिणाम। \\
\textbf{Bioinformatics} (जैव-सूचना विज्ञान) & जैविक data (DNA, प्रोटीन) का कंप्यूटर द्वारा विश्लेषण। \\
\textbf{Catalyst} (उत्प्रेरक) & वह पदार्थ जो अभिक्रिया की दर बढ़ाता है; AI नए उत्प्रेरक खोजने में सहायक। \\
\textbf{Causation} (कारण-सम्बन्ध) & एक घटना का दूसरी घटना का वास्तविक कारण होना (केवल correlation से भिन्न)। \\
\textbf{Citizen Science} (नागरिक विज्ञान) & आम नागरिकों द्वारा वैज्ञानिक data एकत्र करना। \\
\textbf{Classification} (वर्गीकरण) & वस्तुओं या data को श्रेणियों में बाँटना (जैसे अम्ल/क्षार)। \\
\textbf{Cleaning} (सफ़ाई) & data से त्रुटियाँ, दोहराव एवं रिक्त मान हटाना (चरण~2)। \\
\textbf{Collection} (संग्रह) & प्रयोग, सेंसर या सर्वेक्षण से data एकत्र करना (चरण~1)। \\
\textbf{Computational Physics} (संगणनात्मक भौतिकी) & कंप्यूटर-सिमुलेशन एवं data द्वारा भौतिकी का अध्ययन। \\
\textbf{Conclusion} (निष्कर्ष) & परिणाम की वैज्ञानिक व्याख्या एवं सत्यापन (चरण~5)। \\
\textbf{Correlation} (सहसम्बन्ध) & दो राशियों का साथ-साथ बदलना (कारण होना आवश्यक नहीं)। \\
\textbf{Data} (data) & मापे गए मानों या जानकारी का संग्रह। \\
\textbf{Data Science} (data विज्ञान) & data से ज्ञान एवं निर्णय निकालने का विज्ञान। \\
\textbf{Decision} (निर्णय) & उपलब्ध विकल्पों में से सर्वोत्तम चुनना। \\
\textbf{Deep Learning} (गहन अधिगम) & अनेक परतों वाले neural network पर आधारित ML। \\
\textbf{Diversity} (विविधता) & data का सभी प्रकार की स्थितियों का प्रतिनिधित्व करना। \\
\textbf{Error} (त्रुटि) & मापे गए एवं वास्तविक मान का अंतर। \\
\textbf{Fairness} (निष्पक्षता) & AI द्वारा किसी समूह के साथ पक्षपात न होना। \\
\textbf{Feature} (अभिलक्षण) & data का वह गुण जिस पर model आधारित होता है (जैसे height)। \\
\textbf{GC content} (जी-सी अनुपात) & DNA में G एवं C क्षारकों का प्रतिशत; कई जैविक गुणों से जुड़ा। \\
\textbf{General AI} (सामान्य AI) & मनुष्य-जैसी सर्व-कार्य बुद्धि (अभी केवल शोध का विषय)। \\
\textbf{Hallucination} (मतिभ्रम) & AI द्वारा गढ़ा गया झूठा “तथ्य” जो वास्तव में नहीं होता। \\
\textbf{Large Hadron Collider} (वृहत् हैड्रॉन संघट्टक, LHC) & CERN का विशाल कण-त्वरक जो प्रति सेकंड लगभग एक अरब टकराव करता है। \\
\textbf{Law} (नियम) & प्रकृति के व्यवहार का सामान्य एवं परीक्षित कथन। \\
\textbf{Learning} (अधिगम) & उदाहरणों से नियम या pattern सीखना। \\
\textbf{Least Squares Method} (अल्पतम वर्ग विधि) & त्रुटि का वर्ग न्यूनतम कर सर्वोत्तम रेखा खोजना। \\
\textbf{Machine Learning} (मशीन लर्निंग, ML) & data से स्वयं नियम सीखने वाली AI; AI का एक भाग। \\
\textbf{Modelling} (मॉडलन) & pattern से नियम या पूर्वानुमान बनाना (चरण~4)। \\
\textbf{Molecule Design} (अणु-अभिकल्प) & वांछित गुण वाले नए अणु AI द्वारा सुझाना। \\
\textbf{Narrow AI} (सँकरी AI) & केवल एक कार्य में निपुण AI; आज की लगभग सारी AI इसी प्रकार की। \\
\textbf{Neural Network} (तंत्रिका नेटवर्क) & मस्तिष्क की कोशिकाओं से प्रेरित एक ML model। \\
\textbf{Pattern} (pattern) & data में दिखने वाली नियमितता या रुझान। \\
\textbf{Perception} (प्रत्यक्षण) & चित्र, ध्वनि या सेंसर से जानकारी ग्रहण करना। \\
\textbf{Prediction} (पूर्वानुमान) & ज्ञात data से किसी अनदेखे मान का अनुमान। \\
\textbf{Privacy} (निजता) & व्यक्तिगत data की सुरक्षा एवं गोपनीयता। \\
\textbf{Property Prediction} (गुण-पूर्वानुमान) & अणु या पदार्थ के गुण (जैसे विषाक्तता) पहले से आँकना। \\
\textbf{Protein Folding Problem} (प्रोटीन-वलन समस्या) & अमीनो-अम्ल क्रम से प्रोटीन की 3D संरचना का पूर्वानुमान। \\
\textbf{Random Error} (यादृच्छिक त्रुटि) & हर बार बदलती हुई त्रुटि (जैसे प्रतिक्रिया-समय)। \\
\textbf{Reasoning} (तर्क) & उपलब्ध जानकारी से निष्कर्ष निकालना। \\
\textbf{Regression} (समाश्रयण) & data पर सर्वोत्तम रेखा या वक्र फिट करना। \\
\textbf{Reproducibility} (पुनरावृत्ति) & किसी परिणाम का दोबारा दोहराया जा सकना। \\
\textbf{Responsible AI} (उत्तरदायी AI) & पारदर्शिता, निष्पक्षता, जवाबदेही एवं निजता के साथ AI का उपयोग। \\
\textbf{Second Opinion} (द्वितीय राय) & निर्णय की पुष्टि हेतु एक अतिरिक्त राय (जैसे रोग-निदान में AI)। \\
\textbf{SMILES} (स्माइल्स संकेतन) & अणु को एक पंक्ति (string) के रूप में लिखने की विधि। \\
\textbf{Sufficient Quantity} (पर्याप्त मात्रा) & pattern स्पष्ट दिखने योग्य data की मात्रा। \\
\textbf{Systematic Error} (क्रमबद्ध त्रुटि) & हर बार एक ही दिशा में होने वाली त्रुटि। \\
\textbf{Training Data} (प्रशिक्षण data) & model को सिखाने हेतु दिए गए उदाहरण। \\
\textbf{Transparency} (पारदर्शिता) & AI के किसी निर्णय की व्याख्या कर पाने की योग्यता। \\
\textbf{Virtual Screening} (आभासी छँटाई) & कंप्यूटर पर ही लाखों अणुओं को जाँचकर सम्भावित दवाएँ छाँटना। \\
\end{longtable}
